\section*{第 \ref{ch:g6:life-through-seasons} 章——四季中的生命}

\begin{solution}{exo:g6:life-through-seasons:1}
日长（光照）、\omterm{def:g6:exploring-our-environment:conditions}{温度}，以及
降雨和干旱之后土壤或空气的\omterm{def:g6:exploring-our-environment:conditions}{湿度}。
\end{solution}

\begin{solution}{exo:g6:life-through-seasons:2}
落叶——秋季落下全部叶片：橡树、山毛榉、七叶树。
\omterm{def:g6:life-through-seasons:trees}{常绿}——整个
冬季都保持工作中的叶片：松树、冬青。
\end{solution}

\begin{solution}{exo:g6:life-through-seasons:3}
一个微小而完整的芽——微型折叠叶片包裹在
防水鳞片中。它是前一个夏天形成的。
\end{solution}

\begin{solution}{exo:g6:life-through-seasons:4}
\omterm{prop:g6:life-through-seasons:herbs}{一年生植物}整株死亡，以\omterm{def:g2:from-seed-to-plant:seed}{种子}度过冬季。
水仙的地上部分枯萎，而植株
以\omterm{ex:g4:plants-without-seeds:bulbtuber}{鳞茎}的形式在地下等待。
\end{solution}

\begin{solution}{exo:g6:life-through-seasons:5}
因为芽内的幼枝在去年夏天就已经发育完成，
小枝只缺少温暖，而室内的水
比花园里的水早几周提供温暖。
\end{solution}

\begin{solution}{exo:g6:life-through-seasons:6}
\omterm{def:g3:animals-through-seasons:migration}{迁徙}：燕子。\omterm{def:g3:animals-through-seasons:hibernation}{冬眠}：刺猬（或封闭壳口的
蜗牛）。\omterm{def:g3:animals-through-seasons:active}{保持活动}：狐狸、喜鹊、知更鸟。以
幼体阶段越冬——卵、幼虫或蛹：大多数\omterm{prop:g3:sorting-living-things:groups}{昆虫}，
如蚱蜢或蝴蝶。
\end{solution}

\begin{solution}{exo:g6:life-through-seasons:7}
蝴蝶：以蛹的形式藏在压顶石下。蚱蜢：
卵产在土壤中。瓢虫：成虫深深蜷缩在
落叶层中\omterm{prop:g1:healthy-habits:needs}{睡眠}。蜗牛：封闭在岸边的壳中。燕子：
在非洲，以\omterm{def:g3:animals-through-seasons:migration}{迁徙}后的身份生活。
\end{solution}

\begin{solution}{exo:g6:life-through-seasons:8}
这样，四张表之间的任何差异都来自
季节，而不是来自调查本身——
即\cref{met:g6:exploring-our-environment:survey}中的公平试验原则，
只是从空间比较改为时间上的比较。
\end{solution}

\begin{solution}{exo:g6:life-through-seasons:9}
首先，许多“缺席者”以隐蔽的形式存在——
卵、蛹、\omterm{prop:g1:healthy-habits:needs}{睡眠}的成虫、\omterm{ex:g4:plants-without-seeds:bulbtuber}{鳞茎}——
边走边看的调查若不挖掘和翻动，就会漏掉它们。
其次，有些并非消失，而是迁往了别处：\omterm{def:g3:animals-through-seasons:migration}{迁徙动物}
将会返回。一月份的清单低估了留居者，
并把\omterm{def:g3:animals-through-seasons:migration}{迁徙}者误记为已消失。
\end{solution}

\begin{solution}{exo:g6:life-through-seasons:10}
该\omterm{def:g1:plants-around-us:plant}{植物}处于地下储藏状态——仍然活着，
以\omterm{ex:g4:plants-without-seeds:bulbtuber}{鳞茎}的形式等待。三月份\omterm{ex:g4:plants-without-seeds:bulbtuber}{鳞茎}将再次长出叶和
花，若株丛已经增殖，则会更加密集。
\end{solution}

\begin{solution}{exo:g6:life-through-seasons:11}
去了哪里，以什么形式存在？青蛙以休眠方式越冬——
埋在池塘底部的淤泥中，或躲在陆地上潮湿的庇护所里——
在寒冷的月份中处于迟钝的蛰伏状态，
直到春天来临，才被唤回水中繁殖。
\end{solution}

\begin{solution}{exo:g6:life-through-seasons:12}
冬季对工作中的叶片而言，难题在于保住水分：冰冻的
土壤几乎不能为根提供水分，而风和阳光仍在不断
把水分抽出——这正是鼠妇整个冬天面临的干燥难题。
\omterm{def:g6:life-through-seasons:trees}{常绿}\omterm{def:g1:plants-around-us:plant}{植物}叶片的蜡质和针状结构，把损失
降到树木能够承受的程度。宽而薄的叶片无法如此严密地
封闭；落叶树更经济的办法是每年秋天舍弃薄叶，
把春天的希望以
芽的形式储存起来。
\end{solution}

\begin{solution}{pb:g6:life-through-seasons:1}
\textbf{1.} 每根小枝上的芽——冬季储存的养分，
用于下一个春季的新枝，由前一个夏天形成。

\textbf{2.} 三月和四月的画面之间（结合笔记中
日照变长、清晨变暖的情况）。树木
对上升的气温和变长的日照作出反应。

\textbf{3.} \omterm{def:g5:flower-to-fruit:steps}{传粉}\omterm{prop:g3:sorting-living-things:groups}{昆虫}——尤其是蜜蜂——
必须访花，使\omterm{def:g3:flowers-fruits-seeds:parts}{花粉}到达\omterm{def:g3:flowers-fruits-seeds:parts}{雌蕊}，
随后发生\omterm{def:g5:flower-to-fruit:steps}{受精}。七叶树的栗子是由\omterm{def:g5:flower-to-fruit:steps}{受精}的
\omterm{prop:g5:flower-to-fruit:roles}{胚珠}形成的\omterm{def:g2:from-seed-to-plant:seed}{种子}，外面包着多刺的果壳。

\textbf{4.} 落叶。其下一个春天的新枝
以光秃小枝上芽的形式，出现在冬季的画面中。

\textbf{5.} 草坪下方，以\omterm{ex:g4:plants-without-seeds:bulbtuber}{鳞茎}的形式存在——
这是\omterm{prop:g6:life-through-seasons:herbs}{多年生植物}的地下储藏器官。

\textbf{6.} 以\omterm{def:g2:from-seed-to-plant:seed}{种子}的形式留在土壤中：这是\omterm{prop:g6:life-through-seasons:herbs}{一年生
植物}的生存策略——植株本身已在秋天死亡。

\textbf{7.} 燕子：在非洲，以\omterm{def:g3:animals-through-seasons:migration}{迁徙}后的身份
追随飞行的\omterm{prop:g3:sorting-living-things:groups}{昆虫}生活。蜜蜂：在蜂巢中，以
冬季蜂团的形式存在，围绕蜂王保持温暖。

\textbf{8.} \omterm{def:g3:animals-through-seasons:active}{保持活动}。喜鹊必须在最贫乏的几个月里
找到食物——这正是燕子借离开
而避开的难题。

\textbf{9.} 例如：芽\omterm{def:g2:from-seed-to-plant:germination}{萌发}（三月至四月），此时白昼达到
$11$--$13$ 小时，清晨气温升过
\qty{8}{\degreeCelsius}；叶变黄并脱落（十月至十一月）时，
白昼变短，清晨变冷。
树木的各项活动与\omterm{def:g6:exploring-our-environment:conditions}{条件}变化逐帧相对应。

\textbf{10.} 斑块并未死亡：它退入了储藏状态——
以灯泡的形式活着，藏在原处的地下。“消失”总是
值得追问：消失到\emph{哪里}，又以\emph{何种}形式存在？

\textbf{11.} 光秃的枝条上带着芽，附近很可能有一只喜鹊；
地面下，落叶覆盖着光秃的水仙地，
下面藏着灯泡；小路的土壤中有罂粟\omterm{def:g2:from-seed-to-plant:seed}{种子}。确定的是
芽与光秃——它们源于
树木的落叶机制，以及一年的观察证据；喜鹊
则只是很可能，并非确定。

\textbf{12.} 这说明温暖才是缺失的\omterm{def:g6:exploring-our-environment:conditions}{条件}——
二月室内的光照并不比室外更长，但水和房间都是温暖的，
这已足够。公平的对照：
取同一棵树上的两根枝条，同时插在同一扇窗前的水中，
一根放在温暖的房间，一根放在寒冷的门廊——
唯一的差别是\omterm{def:g6:exploring-our-environment:conditions}{温度}；若只有温暖的那根长出\omterm{def:g1:plants-around-us:parts}{叶子}，
结论便确凿无疑。
\end{solution}
